\section{Introduction}
\label{sec:introduction}

Prior work showed that gradient descent on log-sum-exp distance objectives
implements implicit expectation-maximization in a precise local sense: the
gradient with respect to each component distance is the component
responsibility \citep{oursland2025implicitem}. A decoder-free single-layer
encoder tested the unsupervised consequences of that identity and confirmed
the predicted local behavior \citep{oursland2026decoderfree}.

That study also exposed a more restrictive property. The single-layer model
had optimizer dynamics that ordinary neural networks do not generally exhibit:
SGD was insensitive to learning rate over a wide range, Adam offered little
representational advantage, lower loss did not necessarily imply better
features, and convergence time was roughly fixed. These behaviors are
EM-like in a stronger sense than the gradient identity alone. They suggest
that the local responsibility structure can condition the effective parameter
update.

Those optimizer dynamics do not hold for neural networks in general, even
though many standard objectives contain log-sum-exp or softmax structure.
Cross-entropy classifiers, attention layers, contrastive objectives, and
energy-based models all contain exponentiation-and-normalization interfaces,
but they do not usually train with the learning-rate insensitivity of the
single-site implicit-EM model. The question is therefore not whether
responsibility gradients exist. The question is where, under composition, the
optimizer consequences of those gradients are lost.

This paper studies that question in the supervised regime. The testbed is a
minimal supervised network with an explicit intermediate
distance-to-responsibility site. The first layer produces component distances,
a NegLogSoftmin transform creates a calibrated EM interface, and a learned
dense map $W_2$ sends those component coordinates to class logits. A local
auxiliary objective gives the intermediate site the same volume-control
machinery used in the decoder-free model, while the supervised
cross-entropy objective pulls gradients back through the dense head.

The outcome separates two EM properties. Volume control transfers: when
applied at the intermediate site, it eliminates dead units and reduces
redundancy across widths. The optimizer signature does not transfer under
ordinary supervised composition: SGD becomes learning-rate sensitive again,
and adaptive methods recover their advantage. The network contains a valid
responsibility-gradient site, but the deeper supervised model no longer
behaves like the isolated implicit-EM model.

The mechanism is ownership of the update. The local identity creates an
EM-structured path at the intermediate site, but the optimizer behavior at
$W_1$ depends on whether that path controls the effective update. In the
supervised model, the cross-entropy gradient reaching $W_1$ is pulled back
through $W_2$ and added to the local EM gradient. When that supervised
corridor dominates, the update at $W_1$ is governed by a composed, non-EM
contribution even though the local responsibility site still exists.

The resulting boundary has three parts. Implicit EM is \emph{local}: the
responsibility structure exists at LSE and softmax sites, and it extends to
parameters only along paths that preserve component identity. It is
\emph{graded}: the relevant variable is the measured share of EM-structured
gradient in the effective update, not merely whether the computation graph is
connected. It is \emph{structural}: at-site failures such as dead or redundant
components are different from corridor failures caused by learned dense
composition.

The theory formalizes this boundary. In distance coordinates, the LSE Hessian
has spectral norm at most $1/2$. With private affine component parameters,
this gives a parameter-independent curvature bound. When the same site is
trained through a learned dense supervised corridor, the leading
Gauss--Newton term of the cross-entropy Hessian is a conjugated softmax
Hessian whose scale depends on the learned map. The local LSE guarantee no
longer controls the composed update.

The experiments show the same split. Volume-control ablations confirm that
local mixture failures can be repaired inside a supervised model. Optimizer
sweeps show that the single-site signature disappears under ordinary
supervised composition. A causal intervention identifies the mechanism:
blocking the CE corridor to $W_1$ with stop-gradient restores learning-rate
insensitivity, and increasing the EM weight until the local path dominates
restores the same behavior without cutting the graph. Gradient diagnostics
match the intervention. At the standard auxiliary weight, the CE corridor
dominates the weighted local EM path by roughly $67$--$95\times$ late in
training, with no detectable alignment or opposition between the two
directions.

This paper makes the following contributions.
\begin{enumerate}
    \item It identifies the missing condition for EM-like optimizer dynamics
    in composed networks: the responsibility-gradient path must directly or
    dominantly control the relevant parameter update.
    \item It proves a conditioning contrast between private-affine LSE
    objectives and LSE sites composed behind learned dense maps, and
    identifies operations that preserve or destroy the responsibility-gradient
    invariant.
    \item It separates at-site failures from structural corridor failures in
    supervised networks with intermediate EM sites.
    \item It shows that volume control transfers to intermediate supervised EM
    sites, repairing dead or redundant components, while the single-site
    optimizer signature does not survive ordinary dense supervised
    composition.
    \item It restores the conditioning signature by stop-gradient and by EM
    dominance, and uses the $1/2$-saturation diagnostic to measure local
    mixture sharpness at the EM site.
\end{enumerate}
